\section{Check 2: Splits}\label{sec:splits}

\emph{What a random split hides.} Table~\ref{tab:perclass} gives per-class results on CIC-IoT2023. Macro-F1 is 0.867 under the random split and 0.725 under the temporal split. A paired bootstrap puts the difference at $+0.142$ [0.025, 0.287], and almost all of it is one class: BruteForce contributes 78\% of the summed per-class F1 difference. Without BruteForce the difference is $+0.036$. How certain these differences are depends on the unit that is resampled (Fig.~\ref{fig:gapunits}): flows of one host pair are not independent, and coarser units (attacker host, attacker device, capture) are fewer. The full gap excludes zero at every unit but attacker device; the gap without BruteForce excludes zero by host pair and attacker host but not by attacker device or capture. Compared capture by capture, the random split scores higher in 6 of the 9 captures where the scores differ, which is not significant (sign test $p=0.51$). The claim for CIC-IoT2023 is therefore narrow: a random split hides a change of target in one class; beyond it, the gap is small and not robust at the coarser units, which have little power (nine captures, mostly one attack class each).

\begin{figure}[t]
\centering
\includegraphics[width=\columnwidth]{gap_units.pdf}
\caption{Random minus temporal macro-F1 on CIC-IoT2023 (deployed tree), with 95\% paired bootstrap CIs when whole host pairs, attacker hosts, attacker devices or captures are resampled. Macro-F1 is averaged over the classes present in each resample, since coarse resamples often miss a class.}
\label{fig:gapunits}
\end{figure}

\begin{table}[t]
\centering
\caption{Per-class results of the deployed tree on CIC-IoT2023. $n$ and pairs: temporal test flows and host pairs. Recall: temporal, with 95\% host-pair CI; none is given for one pair or for recall 1.000 in every pair, and intervals over at most six pairs are indicative only.}
\label{tab:perclass}
\footnotesize
\setlength{\tabcolsep}{3pt}
\begin{tabular}{lrrcccl}
\hline
 & & & \multicolumn{2}{c}{F1} & \multicolumn{2}{c}{Temporal recall} \\
Class & $n$ & Pairs & random & temp. & & host-pair CI \\
\hline
Benign     &  52,993 & 3,424 & 0.992 & 0.974 & 0.962 & [0.901, 0.990] \\
BruteForce &     605 &     3 & 0.988 & 0.100 & 0.053 & [0.000, 0.842] \\
DDoS-SYN   & 153,402 &     6 & 1.000 & 1.000 & 1.000 & every pair 1.000 \\
DoS-HTTP   &  10,110 &     1 & 0.988 & 0.919 & 0.985 & 1 pair, no interval \\
MITM-ARP   &      25 &    13 & 0.101 & 0.018 & 0.160 & [0.000, 0.375] \\
Mirai      &      10 &     6 & 1.000 & 0.900 & 0.900 & [0.571, 1.000] \\
Recon      &  39,819 &    18 & 1.000 & 1.000 & 1.000 & every pair 1.000 \\
Web        &     192 &     8 & 0.865 & 0.889 & 0.917 & [0.804, 0.941] \\
\hline
Macro-F1   & 257,156 &       & 0.867 & 0.725 & & \\
\hline
\end{tabular}
\end{table}

\emph{A change of attack target.} The brute-force capture shows what the gap is. Training holds only SSH flows from the attacker (.65) to one server on port 22. The test period holds the last 38 SSH flows of that pair and 567 RTSP flows to two cameras on port 554, whose pairs never appear in training. The RTSP flows guess stream paths, not passwords: each short connection sends one \texttt{DESCRIBE} of a path from a dictionary (\texttt{/video.cgi}, \texttt{/live/ch00\_1}, \ldots). The deployed tree detects 32 of the 38 SSH flows and none of the RTSP flows: its recall of 0.053 comes entirely from the SSH tail. A random split puts RTSP flows into training and hides that the model learned one service and method, not ``brute force''. This is a change of target and method within one capture session, not drift, and it rests on three host pairs; we treat it as a case study.

\emph{Boundary-spanning connections.} A temporal cut by flow start does not cut connections: a long connection yields several flows (active timeout 120\,s). 36\% of the temporal test flows share a 5-tuple with a training flow of the same capture, i.e.\ continue a connection that started in training; we call these connections boundary-spanning. Removing them lowers macro-F1 from 0.725 to 0.604, mostly because Mirai is then left without test flows; over the other classes the change is small (0.700 to 0.690). All 10 Mirai test flows are such continuations, so the Mirai test score is not evidence of generalization.

\emph{What flow features cannot see.} MITM-ARP reaches precision 0.01 and recall 0.16 on the temporal split. An ARP-spoofing flow is a victim's ordinary traffic relayed through the attacker, and the MAC address, the only distinguishing field, is not a model input.

\begin{table}[t]
\centering
\caption{IoT-23, deployed tree trained on Zeek labels. 95\% CIs from a flow-level and a host-pair bootstrap. Common: macro-F1 over the four classes present in both halves of both splits. Spanning: share of boundary-spanning test flows.}
\label{tab:iot23split}
\footnotesize
\setlength{\tabcolsep}{2.5pt}
\begin{tabular}{lccccc}
\hline
Split & Macro-F1 & Flow CI & Host-pair CI & Common & Spanning \\
\hline
Temporal & 0.649 & [0.648, 0.650] & [0.613, 0.655] & 0.973 & 4.1\% \\
Random   & 0.863 & [0.764, 0.923] & [0.638, 0.939] & 0.956 & 5.8\% \\
\hline
\end{tabular}
\end{table}

\emph{IoT-23: the gap does not replicate.} On IoT-23 the temporal macro-F1 is again lower than the random one (Table~\ref{tab:iot23split}), but because of missing classes: under the temporal split DDoS has no training flows and FileDownload no test flows, so both score 0 by construction. On the four classes present in both halves the temporal split scores slightly higher, but host pairs are near-singletons here, and with clusters by originating host or by capture the paired interval of the difference includes zero. Boundary-spanning flows concentrate in C\&C; removing them leaves the temporal score unchanged but lowers the random common-class macro-F1 from 0.956 to 0.918. We therefore do not claim that the CIC-IoT2023 gap replicates. What carries over is Check~2 (Table~\ref{tab:checklist}): class presence in both halves, per-class recall with host-pair counts, cluster-unit sensitivity and the boundary-spanning share.
